\section{Finite sets of compact local centres, after McCallum}
\label{sec:irreg-Tate}

McCallum bounds the local zero set of a suitable Coleman integral on
a wild cyclic quotient. We show that one fixed integral contains
the centres of every coherent compact Tate-Selmer datum. The
finite-level reciprocity argument is the step that supplies this
uniformity for integral cohomology classes.

For an abelian variety \(B/\Q_v\), let
\[
 \kappa_{p,v}:B(\Q_v)\longrightarrow H^1(\Q_v,T_pB)
\]
be the inverse-limit Kummer map; the analogous global map is denoted
\(\kappa_p:B(\Q)\to H^1(\Q,T_pB)\).

\begin{definition}\label{def:irreg-data}
A \emph{compact \(p\)-Selmer datum on \(F_p\)} consists of
\[
 t\in H^1(\Q,T_pJ_p),\qquad
 R_v\in F_p(\Q_v)\quad(v<\infty),
\]
such that
\begin{equation}
 \loc_v(t)=\kappa_{p,v}\bigl(j_p(R_v)\bigr)
 \quad\text{for every finite }v.
 \label{eq:irreg-data}
\end{equation}
\end{definition}

This definition uses integral Tate-module cohomology. It does not
mean an arbitrary Selmer class on the Jacobian without curve-valued
local representatives. It imposes no common finite set of
ramification primes on \(t\).

A quotient section with full local representatives supplies such a
datum whenever it retains the proper geometrically abelian pro-\(p\)
quotient. Indeed, the base point \(Q_0\) splits that quotient, whose
geometric group is \(T_pJ_p\); sections then give classes in
\(H^1(\Q,T_pJ_p)\), and point sections give Kummer classes.
The geometric identification follows at each finite level from
the Kummer sequence on \(F_{p,\overline\Q}\), the identification of
its degree-zero Picard variety with \(J_p\), and the principal
polarization of \(J_p\).

Conversely, reducing a datum modulo \(p\) gives a section of
\(\Pi_p^{\ab,p}\) with the prescribed proper local point
representatives. Thus the finite first-case criteria above apply
to every datum in Definition~\ref{def:irreg-data}.

For a principally polarized abelian variety \(A/\Q\), write
\begin{equation}
 H^1_{\{p\}}(\Q,A)
 =
 \ker\left(H^1(\Q,A)\longrightarrow
                  \prod_{v\neq p}H^1(\Q_v,A)\right).
 \label{eq:irreg-supported}
\end{equation}
The product includes the infinite place. This is the group of
Weil--Chatelet classes supported at \(p\), not the usual
everywhere locally trivial Tate--Shafarevich group.

\begin{lemma}\label{lem:irreg-reciprocity}
Let \(A/\Q\) be principally polarized, and let
\(u\in H^1(\Q,T_pA)\) have localizations
\(\kappa_{p,v}(P_v)\), with \(P_v\in A(\Q_v)\) at every finite place.
For every
\(d\in H^1_{\{p\}}(\Q,A)[p^n]\),
\begin{equation}
 \langle P_p,d_p\rangle_p=0
 \label{eq:irreg-supported-pairing}
\end{equation}
under local Tate duality.
\end{lemma}

\begin{proof}
Use the principal polarization to identify \(A\) with its dual.
Let \(u_n\in H^1(\Q,A[p^n])\) be the finite image of \(u\), and choose
\(b_n\in H^1(\Q,A[p^n])\) mapping to \(d\). Such a lift exists by
the Kummer exact sequence. The Weil pairing gives a global class
\[
 u_n\smile b_n\in H^2(\Q,\mu_{p^n})
                 \subseteq\operatorname{Br}(\Q).
\]
Compatibility of finite Kummer duality with the local Tate pairing
gives
\[
 \inv_v\bigl((u_n\smile b_n)_v\bigr)
     =\langle P_v,d_v\rangle_v.
\]
A consistent choice of the order of the dual factors fixes the sign;
only its vanishing is used. This compatibility is the finite-level
diagram in the proof of
\cite[I, Theorem 3.2, pp.~41--43]{Milne2006ADT};
the perfect local pairing is
\cite[I, Corollary 3.4, pp.~43--44]{Milne2006ADT}.
It also shows that changing the lift \(b_n\) by a Kummer class does
not change the displayed value.

For \(v\neq p\), the value is zero because \(d_v=0\). At the real
place the \(p\)-primary term is zero, since \(p\) is odd.
Global reciprocity now makes the remaining value at \(p\) zero;
equivalently one may use the orthogonality in
\cite[I, Theorem 4.10(b), pp.~56--57]{Milne2006ADT}.
This argument is made at one finite level at a time and needs no
uniform finite-ramification hypothesis on \(u\).
\end{proof}

\begin{lemma}\label{lem:irreg-tower-functional}
Suppose a compatible tower
\begin{equation}
 d=(d_n)_{n\geq1},\qquad
 d_n\in H^1_{\{p\}}(\Q,A)[p^n],\qquad
 p\,d_{n+1}=d_n,
 \label{eq:irreg-tower}
\end{equation}
has nonzero localization at \(p\). It determines a nonzero
continuous homomorphism
\begin{equation}
 \ell_d:A(\Q_p)\longrightarrow\Z_p.
 \label{eq:irreg-integral-functional}
\end{equation}
Its kernel contains all torsion and is saturated for multiplication
by every nonzero integer. There is a nonzero invariant differential
\(\omega_d\) on \(A/\Q_p\) such that
\begin{equation}
 \ell_d(P)=\omega_d\bigl(\log_A(P)\bigr)
 \quad\text{for all }P\in A(\Q_p).
 \label{eq:irreg-log-functional}
\end{equation}
Every collection \(P_v\) occurring in
Lemma~\ref{lem:irreg-reciprocity} satisfies \(\ell_d(P_p)=0\).
\end{lemma}

\begin{proof}
For each \(n\), pair \(P\) with \(d_{n,p}\) and identify
\[
 p^{-n}\Z/\Z\xrightarrow{\sim}\Z/p^n\Z,
 \qquad a/p^n\longmapsto a\bmod p^n.
\]
Compatibility of the tower gives a compatible system of residues,
hence \(\ell_d(P)\in\Z_p\). Every finite coordinate is continuous.
Perfect local Tate duality shows that this homomorphism is nonzero:
otherwise every \(d_{n,p}\) would vanish.

The target \(\Z_p\) is torsion-free, proving the assertions about
torsion and saturation. Choose an open subgroup \(H\subset A(\Q_p)\)
on which the logarithm identifies \(H\) with a full
\(\Z_p\)-lattice in \(\Lie(A)\). Such a subgroup exists by
\cite[I, Lemma 3.3, p.~41]{Milne2006ADT}.
The restriction of \(\ell_d\) becomes a continuous additive map on
that lattice, hence a \(\Z_p\)-linear map. Extend it uniquely to a
\(\Q_p\)-linear functional on \(\Lie(A)\).
This extension is nonzero: a homomorphism from the finite quotient
\(A(\Q_p)/H\) to \(\Z_p\) must vanish.

For arbitrary \(P\), choose a positive integer \(N\) with \(NP\in H\).
Both the logarithm and \(\ell_d\) commute with multiplication by
\(N\); division by \(N\) in \(\Q_p\) proves
\eqref{eq:irreg-log-functional}. The functional is the invariant
differential \(\omega_d\). Finally,
Lemma~\ref{lem:irreg-reciprocity} makes every finite coordinate of
\(\ell_d(P_p)\) zero.
\end{proof}

\subsection{Wild cyclic quotients}
\label{sec:irreg-wild}

Let \(a,b,c\) be integers with \(a+b+c=0\) and
\(p\nmid abc\). Choose \(a,b>0\) in their residue classes and set
\(c=-a-b\). Let \(C_{a,b,c}\) be the smooth projective curve with
affine equation
\begin{equation}
 y^p=(-1)^c x^a(1-x)^b.
 \label{eq:irreg-cyclic}
\end{equation}
It has genus \((p-1)/2\), by Riemann--Hurwitz for its degree-\(p\)
map to the \(x\)-line, totally ramified over \(0,1,\infty\).
Denote the unique points over these three values by
\(O_0,O_1,O_\infty\). They are rational. Put
\[
 A=\Jac(C_{a,b,c}),\qquad j(R)=[R-O_0].
\]
There is a morphism
\begin{equation}
 f:F_p\longrightarrow C_{a,b,c},\qquad
 x=X^pZ^{-p},\qquad y=X^aY^b(-Z)^c.
 \label{eq:irreg-quotient-map}
\end{equation}
The rational functions define the morphism on the smooth projective
curves by extension from their common affine domain.

\begin{lemma}\label{lem:irreg-quotient-free}
The morphism \(f\) is the quotient by the free action
\begin{equation}
 \mu_p\times F_p\longrightarrow F_p,\qquad
 \zeta\cdot(X:Y:Z)=(\zeta^bX:\zeta^{-a}Y:Z).
 \label{eq:irreg-action}
\end{equation}
In particular, \(f:F_p(\Q_p)\to C_{a,b,c}(\Q_p)\) is injective.
Moreover \(j(O_0),j(O_1),j(O_\infty)\) are torsion points.
\end{lemma}

\begin{proof}
The functions in \eqref{eq:irreg-quotient-map} are invariant under
\eqref{eq:irreg-action}. Over the \(x\)-line, the function field of
\(F_p\) has degree \(p^2\), whereas
\(\Q(x,y)/\Q(x)\) has degree \(p\), since the valuation \(a\) at
\(x=0\) is not divisible by \(p\). These invariants therefore give
the quotient field, and hence the quotient curve.

At a point with all coordinates nonzero, a projective stabilizer
must have \(\zeta^b=\zeta^{-a}=1\), hence \(\zeta=1\).
At \(X=0\) the stabilizer condition is \(\zeta^{-a}=1\);
at \(Y=0\) it is \(\zeta^b=1\); and at \(Z=0\) it is
\(\zeta^{a+b}=1\). Each again forces \(\zeta=1\), using
\(p\nmid abc\). Thus the action is free in every stratum and the
proper quotient is a \(\mu_p\)-torsor.

If two \(\Q_p\)-points have the same image, the unique geometric
element carrying one to the other is Galois-fixed, hence lies in
\(\mu_p(\Q_p)=\{1\}\). This proves injectivity, including at the
boundary.

Finally,
\[
 \divisor(x)=p(O_0-O_\infty),\qquad
 \divisor(1-x)=p(O_1-O_\infty).
\]
Their classes show that each of the three images under \(j\) is
killed by \(p\).
\end{proof}

Following \cite[pp.~395--398]{McCallum2001Brauer}, set
\[
 \gamma=a^a b^b c^c\in\Z_{(p)}^\times.
\]
The quotient is called \emph{wild} if
\begin{equation}
 \gamma^p\not\equiv\gamma\pmod{p^2}.
 \label{eq:irreg-wild}
\end{equation}

\begin{lemma}\label{lem:irreg-wild-exists}
For every prime \(p\geq5\), at least one of the quotients
\(C_{1,k,-k-1}\), \(1\leq k\leq p-2\), is wild.
\end{lemma}

\begin{proof}
For \(u\in\Z_{(p)}^\times\), define its Fermat quotient by
\[
 q_p(u)=\frac{u^{p-1}-1}{p}\pmod p.
\]
Multiplication modulo \(p^2\) shows that this is a homomorphism
from \(\Z_{(p)}^\times\) to the additive group \(\F_p\).
For \(\gamma_k=k^k(-k-1)^{-k-1}\), one has
\[
 q_p(\gamma_k)=kq_p(k)-(k+1)q_p(k+1).
\]
Here \(q_p(-1)=0\). Also
\((p-1)^{p-1}\equiv1+p\pmod{p^2}\), so \(q_p(p-1)=1\).
The telescoping sum is therefore
\[
 \sum_{k=1}^{p-2}q_p(\gamma_k)
   =q_p(1)-(p-1)q_p(p-1)=1\quad\text{in }\F_p.
\]
At least one summand is nonzero, which is exactly
\eqref{eq:irreg-wild}.
\end{proof}

We use two precise results of McCallum. In the following statements,
the class rank is the full \(p\)-class rank, not only its minus part:
\[
 r=\dim_{\F_p}\Cl(K)/p\Cl(K).
\]

\begin{proposition}[McCallum]\label{prop:irreg-McCallum}
For \(C=C_{a,b,c}\) and \(A=\Jac(C)\) as above:
\begin{enumerate}
\item If \(r<(p+3)/8\), there is a compatible tower
\[
 d_n\in H^1_{\{p\}}(\Q,A)[p^n],\qquad p\,d_{n+1}=d_n,
\]
whose localization at \(p\) is nonzero.
\item If \(C\) is wild, there is an affinoid
\(\mathcal X\subset C(\C_p)\) containing \(C(\Q_p)\) such that,
for every nonzero holomorphic differential \(\eta\) on \(C\),
the Coleman integral
\[
 I_\eta(R)=\int_{O_0}^R\eta
\]
has at most \(p\) zeros in \(\mathcal X\).
\end{enumerate}
\end{proposition}

These are \cite[Propositions 2 and 3, pp.~397--398]{McCallum2001Brauer}.
The construction of the tower is proved there in Section 3,
pp.~399--403; the analytic bound is proved in Section 4,
pp.~403--406. The affine region is the tube of the nonsingular
locus of the indicated regular model. Every \(\Q_p\)-point belongs
to it in the wild case. The zero bound applies to any nonzero
holomorphic differential; it does not require a complex
multiplication eigendifferential. The relation of the Coleman
integral to the abelian logarithm is given in the same paper,
Section 2, pp.~397--398.

\subsection{One fixed finite set of local centres}
\label{sec:irreg-fixed}

\begin{theorem}\label{thm:irreg-fixed}
Suppose
\begin{equation}
 r=\dim_{\F_p}\Cl(\Q(\zeta_p))/p\Cl(\Q(\zeta_p))
       <\frac{p+3}{8}.
 \label{eq:irreg-Mc-bound}
\end{equation}
There is one finite set \(S_p\subset F_p(\Q_p)\), independent of
the compact \(p\)-Selmer datum and of all its chosen local
representatives, such that
\begin{equation}
 \{Q_0,Q_1,Q_\infty\}\subseteq S_p,\qquad
 |S_p|\leq p,
 \label{eq:irreg-fixed-size}
\end{equation}
and every datum in Definition~\ref{def:irreg-data} has
\(R_p\in S_p\).

More precisely, one may choose any wild quotient \(C_{a,b,c}\);
a single supported tower for its Jacobian, fixed independently of
the datum, determines such a set.
\end{theorem}

\begin{proof}
Choose a wild quotient, which exists by
Lemma~\ref{lem:irreg-wild-exists}, and let \(A=\Jac(C_{a,b,c})\).
By Proposition~\ref{prop:irreg-McCallum}(1), choose once and for all
a supported tower \(d=(d_n)\) with nonzero \(p\)-localization.
Lemma~\ref{lem:irreg-tower-functional} gives the corresponding
nonzero invariant differential \(\omega_d\) on \(A\).
Put
\[
 \eta=j^*\omega_d.
\]
This is nonzero because the Abel--Jacobi map identifies invariant
differentials on the Jacobian with holomorphic differentials on
the curve. Functoriality of the abelian logarithm and Coleman
integration gives
\begin{equation}
 I_\eta(R)=\omega_d\bigl(\log_A(j(R))\bigr)
           =\ell_d(j(R)).
 \label{eq:irreg-fixed-integral}
\end{equation}

Let \(t,(R_v)\) be any compact \(p\)-Selmer datum on \(F_p\).
The induced homomorphism \(f_*:J_p\to A\), with the specified base
points, satisfies \(f_*j_p=jf\). The image \(u=f_*t\) therefore has
localizations
\[
 \loc_v(u)=\kappa_{p,v}(j(f(R_v))).
\]
Lemma~\ref{lem:irreg-reciprocity}, applied to every \(d_n\), and
\eqref{eq:irreg-fixed-integral} imply
\[
 I_\eta(f(R_p))=0.
\]
All choices of \(d,\omega_d,\eta\) preceded the arbitrary datum.

Define
\begin{equation}
 S_p=
 \{R\in F_p(\Q_p): I_\eta(f(R))=0\}.
 \label{eq:irreg-fixed-set}
\end{equation}
The image of \(F_p(\Q_p)\) lies in \(\mathcal X\), where
Proposition~\ref{prop:irreg-McCallum}(2) bounds the zeros of
\(I_\eta\) by \(p\). Lemma~\ref{lem:irreg-quotient-free} makes \(f\)
injective on \(\Q_p\)-points, so \(|S_p|\leq p\).

The three rational coordinate cusps map to \(O_0,O_1,O_\infty\).
Their images in \(A\) are torsion, and \(\ell_d\) kills torsion.
Thus all three belong to \(S_p\), proving \eqref{eq:irreg-fixed-size}.
\end{proof}

\begin{remark}\label{rem:irreg-uniform}
The set \(S_p\) need not be canonical: it depends on a fixed quotient
and a fixed supported tower. Its defining choices are independent
of the Selmer datum, which is the uniformity asserted in the
theorem. Neither a first-case hypothesis nor finiteness of a
Tate--Shafarevich group is used for the bound.
\end{remark}

\begin{corollary}\label{cor:irreg-fixed-discs}
In addition to \eqref{eq:irreg-Mc-bound}, suppose either
\eqref{eq:irreg-rank-bound} or the class-component hypothesis of
Corollary~\ref{cor:irreg-cubic}. Then \(S_p\) may be replaced by its
intersection with the three coordinate cusp residue discs while
retaining all conclusions of Theorem~\ref{thm:irreg-fixed}.
\end{corollary}

\begin{proof}
Reduction of every compact datum modulo \(p\) has the full proper
local representatives required for the first-case arguments.
Those arguments put every allowed centre in the three discs.
The intersection still contains the three cusps.
\end{proof}

\subsection{The primes 37, 59, and 67}
\label{sec:irreg-examples}

Masley's theorem gives class number one for every real abelian
field of prime-power conductor less than \(70\)
\cite[Theorem 3.2, p.~310]{Masley1978}. Newman's factored table gives
the following relative class numbers
\cite[Table 2, p.~218]{Newman1970}:
\begin{equation}
 \begin{array}{c|c|c}
 p&h^+(\Q(\zeta_p))&h^-(\Q(\zeta_p))\\ \hline
 37&1&37\\
 59&1&3\cdot59\cdot233\\
 67&1&67\cdot12739.
 \end{array}
 \label{eq:irreg-class-table}
\end{equation}
The cofactors at \(59\) and \(67\) are nonzero modulo the corresponding
prime:
\[
 699\equiv50\pmod{59},\qquad
 12739\equiv9\pmod{67}.
\]
Thus each \(p\)-Sylow class group has order \(p\). The norm to the
trivial real class group, followed by restriction, is \(1+\overline{\phantom{x}}\);
it follows that the \(p\)-class group is entirely minus. Therefore
\begin{equation}
 r=r^-=1\qquad(p=37,59,67).
 \label{eq:irreg-ranks-three}
\end{equation}
The full class groups at \(59\) and \(67\) have the additional factors
shown in \eqref{eq:irreg-class-table}.

For each of these primes one may take the same quotient
\begin{equation}
 C_p:\ y^p=x(1-x),\qquad
 f_p(X:Y:Z)=(X^p/Z^p,XY/Z^2),
 \label{eq:irreg-111-quotient}
\end{equation}
corresponding to \((a,b,c)=(1,1,-2)\). Here \(\gamma=1/4\).
Exact modular exponentiation gives
\begin{equation}
 \begin{array}{c|c|c}
 p&p^2&4^{p-1}\bmod p^2\\ \hline
 37&1369&75=1+2p\\
 59&3481&945=1+16p\\
 67&4489&1341=1+20p.
 \end{array}
 \label{eq:irreg-wild-table}
\end{equation}
Hence all three quotients are wild.

\begin{corollary}\label{cor:irreg-37-59-67}
For \(p=37,59,67\), there is one fixed set \(S_p\subset F_p(\Q_p)\)
containing the three rational cusps, contained in their three
coordinate residue discs, and satisfying
\[
 3\leq |S_p|\leq p,
\]
which contains the \(p\)-local centre of every compact
\(p\)-Selmer datum on \(F_p\).

For \(p=37\), one may use \(C_{37}:y^{37}=x(1-x)\), of genus \(18\),
to construct \(S_{37}\). In particular, the same \(S_{37}\) works for
every quotient section retaining the proper compact \(T_{37}J_{37}\)
projection and having the stipulated full local representatives.
\end{corollary}

\begin{proof}
Equation \eqref{eq:irreg-ranks-three} gives
\(r<(p+3)/8\), and the rank first-case criterion becomes \(10<p\).
Apply Theorem~\ref{thm:irreg-fixed} and
Corollary~\ref{cor:irreg-fixed-discs}, using the wild quotients in
\eqref{eq:irreg-111-quotient}. Their genus is \((p-1)/2\).
\end{proof}

The cardinality conclusions are upper bounds. They do not assert
\(S_p=\{Q_0,Q_1,Q_\infty\}\).

\subsubsection{Higher exponents}

\begin{corollary}\label{cor:irreg-higher-fixed}
Let \(n=pm\) be odd and assume \eqref{eq:irreg-Mc-bound}. Put
\(d=\gcd(m,p-1)\). There is one fixed finite set
\(S_{n,p}\subset F_n(\Q_p)\) containing every \(\Q_p\)-rational
coordinate boundary point, such that
\begin{equation}
 |S_{n,p}|\leq(p-3)d^2+3d.
 \label{eq:irreg-higher-bound}
\end{equation}
It contains the \(p\)-local centre of every datum consisting of
\[
 t\in H^1(\Q,T_pJ_n),\qquad R_v\in F_n(\Q_v),
 \qquad
 \loc_v(t)=\kappa_{p,v}([R_v-Q_0])
 \quad(v<\infty).
\]
If \(d=1\), the bound is \(p\). In particular \(d=1\) whenever
\(p\) is the least prime divisor of \(n\).
\end{corollary}

\begin{proof}
Let
\[
 \phi_m:F_n\longrightarrow F_p,\qquad
 (X:Y:Z)\longmapsto(X^m:Y^m:Z^m),
\]
and take \(S_{n,p}=\phi_m^{-1}(S_p)\) on \(\Q_p\)-points.
Functoriality on Jacobians pushes the given global Tate class and
the same local representatives to a compact \(p\)-Selmer datum on
\(F_p\). Hence every permitted centre belongs to this inverse image.

The roots of unity in \(\Q_p\) have order dividing \(p-1\), so
\(|\mu_m(\Q_p)|=d\). Above an ordinary point, normalize the nonzero
third coordinate to \(1\). Each of the other two coordinates has
at most \(d\) possible \(m\)-th roots, giving at most \(d^2\) points
in the fibre. Any such choices satisfy the exponent-\(n\) equation
automatically. Above \(Q_0\) or \(Q_1\), only one root choice remains,
giving exactly \(d\) points. Above \(Q_\infty\), the equation for
the coordinate ratio is \(u^m=-1\); since \(m\) is odd, its solutions
are \(-\mu_m(\Q_p)\), again exactly \(d\).

Every rational coordinate boundary point of \(F_n\) maps to one of
the three rational coordinate cusps of \(F_p\). Since those three
belong to \(S_p\), the inverse image contains them all, and
\[
 |S_{n,p}|
 \leq(|S_p|-3)d^2+3d
 \leq(p-3)d^2+3d.
\]
If \(p\) is the least prime divisor of \(n\), every prime dividing
\(m\) is at least \(p\), whereas every prime dividing \(p-1\) is
smaller than \(p\). Thus \(\gcd(m,p-1)=1\).
\end{proof}


The fixed-set theorem and its higher-exponent consequence prove
Theorem~\ref{thm:finite-local-centres-intro}.
